% Generated draft: requires mathematical and source review.
\documentclass[11pt]{article}
\usepackage[margin=25mm]{geometry}
\usepackage{fontspec}
\setmainfont{Latin Modern Roman}
\usepackage{amsmath,amssymb,mathtools,unicode-math}
\setmathfont{Latin Modern Math}
\usepackage{enumitem,needspace,xurl,hyperref}
\setlength{\parindent}{0pt}
\setlength{\parskip}{5pt}
\setlength{\emergencystretch}{3em}
\newcommand{\N}{\mathbb{N}}
\newcommand{\Z}{\mathbb{Z}}
\newcommand{\Q}{\mathbb{Q}}
\newcommand{\R}{\mathbb{R}}
\newcommand{\C}{\mathbb{C}}
\newcommand{\F}{\mathbb{F}}
\newcommand{\A}{\mathbb{A}}
\newcommand{\PP}{\mathbb{P}}
\newcommand{\E}{\mathbb{E}}
\newcommand{\eps}{\varepsilon}
\newcommand{\dd}{\,\mathrm{d}}
\begin{document}
\section*{UnsolvedMath ID: 11100055}
\textbf{Source problem code:} catalogue entry
\subsection{Definitions and mathematical statement}
The source proposes, explicitly as Hovey's expectation, a ``2-Quillen equivalence'' between a 2-category of model categories and one of simplicial model categories. A simplicial model category is enriched, tensored and cotensored over simplicial sets with the compatible model-category axiom. The 1-morphisms, 2-morphisms, weak equivalences and meaning of 2-Quillen equivalence are not defined in this excerpt; the preceding two problems containing the intended context are absent.

The concrete task is to understand Dugger's simplicial-replacement construction functorially. For an admissible model category \(\mathcal C\), it provides a simplicial model category \(R\mathcal C\) with an appropriate Quillen equivalence or zigzag to \(\mathcal C\). The source asks whether these replacements extend to a 2-functor \(R\), with compatible assignments to the selected Quillen morphisms and natural transformations and coherence for identities/composition. Strict 2-functoriality, pseudofunctoriality and a construction only after localization are distinct possibilities; none is selected without the missing definitions.

It further asks whether the preserved structure makes an equivalence of the intended homotopy 2-categories evident. In a bicategorical formulation this would require essential surjectivity on objects and equivalences on the appropriate derived hom-categories, compatibly with composition. Objectwise existence of Quillen-equivalent simplicial replacements alone proves no such local mapping-category statement.

Hovey reports Dugger's result with unspecified hypotheses, tentatively saying at least cofibrant generation. The summary formulates modern results for combinatorial or otherwise suitable categories, and invokes higher-category presentations. A combinatorial model category is cofibrantly generated with locally presentable underlying category; it is not synonymous with every cofibrantly generated model category. Those domain restrictions and the replacement's choices must be fixed before claiming a functor on all model categories. The summary's higher-categorical interpretation is related progress, not automatically a strict 2-functor or the literal undefined 2-Quillen equivalence.
\subsection{Short English statement}
Make Dugger's simplicial replacement coherent on morphisms and transformations, and determine whether it yields the intended homotopy-2-categorical equivalence. The domain hypotheses and higher-categorical definitions are part of the missing formulation.
\subsection{Sources}
\begin{itemize}
\item Local UnsolvedMath registry, permanent ID \(11100055\), exact source code \texttt{catalogue entry}.
\item Mark Hovey, \emph{Hovey's Algebraic Topology Problem List}, Model categories, item 9, surviving mirror: \url{https://www-users.cse.umn.edu/~tlawson/hovey/}. Registry access 2026-07-29; html-ordered-list extraction.
\item Dan Dugger is credited by the source; the summary cites \emph{Universal homotopy theories} (2001) for related replacement/presentation results, without the exact theorem and hypotheses. The preceding two source problems and the truncated modern higher-category discussion are not supplied.
\end{itemize}
\end{document}
